\subsection{Geometry-level sensitivity of completed comparisons}
The mean reduction alone does not describe how broadly improvements occur.
We therefore reconstructed each geometry effect directly from the audited
seed costs of the two completed fresh 60-second studies. Rational arithmetic
was used to check the reported seed-mean effects before rounding. Training
and reused validation were analyzed separately; their different geometry
sets and selection histories preclude pooling them into one confirmation.

\begin{table}[htbp]
\centering
\small
\begin{tabular}{lrrrr}
\toprule
Study & $n$ & Mean (\%) & Median (\%) & LOO mean range (\%) \\
\midrule
TRAIN & 9 & 1.886 & 2.404 & [1.609, 2.134] \\
Reused validation & 7 & 0.907 & 0.000 & [0.430, 1.084] \\
\bottomrule
\end{tabular}
\caption{Descriptive sensitivity of completed 60-second comparisons. The leave-one-geometry-out (LOO) range recomputes the equally weighted mean after each possible single omission. It is not a confidence interval or an alternative primary endpoint. All original geometries remain in the reported primary means.}
\label{tab:geometry-sensitivity}
\end{table}


For the nine training geometries, the median reduction was 2.404\%, compared
with a mean of 1.886\%. Recomputing the mean after each possible single
geometry omission gave values from 1.609\% to 2.134\%
(Table~\ref{tab:geometry-sensitivity}). Thus, no single training geometry
accounted for the entire positive average. This observation concerns the
selected training set and does not establish performance on new geometries.

For the seven reused validation geometries, the median effect was zero,
despite the positive mean of 0.907\%. The single-omission means ranged from
0.430\% to 1.084\%. The breadth of improvement remained limited: four
geometry means tied, one deteriorated and only two improved. As a
retrospective concentration diagnostic, the five geometries other than
pr264 and pr299 had a mean reduction of $-0.030\%$. This is not an alternative
primary comparison: both positive cases remain in the seven-case result.
It explains why the positive overall mean can coexist with failure of the
original requirement for at least three positive geometry means.

The omission ranges are descriptive sensitivity measures, not confidence
intervals. Seed repetitions characterize stochastic variation conditional
on each geometry; they are not additional independent benchmark geometries.
Because several candidate methods had already been examined and validation
geometries had been reused, no confirmatory hypothesis test is attached to
these post hoc diagnostics. They neither change a screening threshold nor
justify stopping the incomplete longer-horizon collection early.

\subsection{Sensitivity of the completed 300-second screen}
\label{sec:awake300-sensitivity}
The same descriptive omission calculation was repeated for both prespecified
300-second variants, using exact rational reconstructions of the 81 audited
endpoint costs. For the 60-second initialization, single-omission means ranged
from 0.840\% to 1.112\%; no individual geometry accounted for the positive
average. This is a statement about the selected training set, not a confidence
interval or evidence of generalization.

For the 10-second initialization, the range was $-1.683\%$ to 0.996\%.
The positive endpoint arose when rd400 was omitted, illustrating concentration
of the observed deterioration. That geometry remains in the primary mean of
$-1.279\%$ and the original screen still fails. Neither the omission nor the
sensitivity range supplies a replacement selection rule. Paired seeds within
one geometry were not treated as independent instances.
